\chapter{IL NORD DI OVUNQUE}\label{il-nord-di-ovunque}

\hfill\break

Un'ora e mezza prima di attraversare l'Arco, un'ora dopo il buio,
incrociammo En alla mensa dell'equipaggio. Uno dei membri gli aveva dato
un foglio di carta marroncina e qualche pastello per tenerlo occupato.

Sembrò contento di vederci. Era preoccupato per il transito, disse.
Spinse gli occhiali in su, sopra il naso - facendo una smorfia quando
con il pollice sfiorò la sbucciatura che gli aveva lasciato Jala sulla
guancia - e mi chiese come sarebbe stato.

«Non lo so» risposi. «Non ho mai fatto il transito.»

«Ce ne accorgeremo quando succede?»

«Secondo l'equipaggio, il cielo diventa un po' strano. E non appena
accade, quando siamo in bilico tra il vecchio e il nuovo mondo, l'ago
della bussola gira al contrario, il Nord diventa Sud. E sul ponte suona
la sirena della nave. Te ne accorgerai.»

«Fare tanta strada» mormorò En. «In poco tempo.»

Era senza dubbio vero. Prima di cadere dall'orbita, l'Arco - almeno il
nostro lato - era stato trascinato fisicamente attraverso lo spazio
interstellare a una velocità probabilmente di poco inferiore a quella
della luce. Ma gli Ipotetici avevano avuto eoni di tempo dello Spin per
trascinarlo. Avrebbero plausibilmente potuto colmare qualunque distanza
minore di tre miliardi di anni luce. E persino una piccola frazione di
quella sarebbe stata una distanza stordente e a malapena comprensibile.

«Viene da chiedersi,» brontolò Diane, «perché si siano presi tanto
disturbo.»

«Secondo Jason\ldots»

«Lo so. Gli Ipotetici vogliono proteggerci dall'estinzione, in modo che
ci rendiamo più complessi. Ma la domanda nasce spontanea. Perché lo
vogliono? Cosa si aspettano?»

En ignorò il nostro filosofeggiare. «E dopo che avremo
attraversato\ldots»

«Dopo quello,» gli spiegai, «ci sarà un giorno di viaggio prima di
arrivare a Port Magellan.»

Sorrise a quella prospettiva.

Scambiai uno sguardo con Diane. Si era presentata a En due giorni prima
ed erano già amici. Gli aveva letto qualcosa da un libro di storie di
bambini inglesi preso nella biblioteca della nave. (Gli aveva pure
citato Housman: ``\emph{All'infante non è occorso}\ldots'' «Non mi piace
quella» aveva detto En.)

Ci mostrò i suoi disegni, immagini di animali che doveva avere visto nei
filmati girati nelle pianure di Equatoria, bestie dal collo lungo con
occhi pensosi e pellicce tigrate.

«Sono bellissimi» si complimentò Diane.

En annuì con una certa solennità. Lo lasciammo al suo lavoro e ci
dirigemmo sul ponte.

\separatorefiori

Il cielo notturno era limpido e l'apice dell'Arco, sopra di noi,
rifletteva un ultimo barlume di luce. Non mostrava alcuna curvatura. Da
quell'angolo era una pura linea euclidea, un numero elementare (1) o una
lettera (I).

Rimanemmo accanto al corrimano, il più vicino possibile alla prua della
nave. Il vento ci strattonava la giacca e i capelli. Le bandiere della
nave schioccavano vivaci e un mare inquieto restituiva immagini spezzate
delle luci accese della nave.

«Ce l'hai?» chiese Diane.

Si riferiva alla minuscola fiala che conteneva un campione delle ceneri
di Jason. Avevamo pianificato questa cerimonia - se si poteva chiamare
cerimonia - molto prima di lasciare Montreal. Jason non aveva mai
riposto troppa fiducia nelle commemorazioni, ma credo che avrebbe
approvato questa. «Proprio qui.» Tirai fuori dalla tasca del gilet il
piccolo tubo di ceramica e lo tenni stretto nella mano sinistra.

«Mi manca» disse Diane. «Mi manca continuamente.» Si insinuò nella mia
spalla, la abbracciai. «Vorrei averlo conosciuto come Quarto. Ma
suppongo che non lo avesse cambiato molto\ldots»

«Già, non lo aveva cambiato molto.»

«In un certo senso Jase è \emph{sempre} stato un Quarto.»

Mentre ci avvicinavamo al momento del transito, le stelle sembravano
affievolirsi, come se una diafana presenza avesse avvolto la nave. Aprii
la fiala con le ceneri di Jason. Diane mise la mano libera sulla mia.

Il vento cambiò direzione all'improvviso, la temperatura si abbassò di
un grado o due.

«A volte,» bisbigliò Diane, «quando penso agli Ipotetici, ho paura
che\ldots»

«Che cosa?»

«Che siamo il loro vitello rosso. O ciò che Jason sperava sarebbero
stati i marziani. Che si aspettino che noi li salviamo da qualcosa.
Qualcosa di cui \emph{loro} hanno paura.»

Forse è così. ``Ma poi,'' pensai, ``faremo ciò che la vita fa di
solito\ldots'' sfuggire alle aspettative.

Sentii che un brivido l'attraversava. Sopra di noi la linea dell'Arco si
stava allontanando. Una nebbiolina si poggiò sul mare. Però non era una
nebbia normale. Non era un fenomeno meteorologico.

L'ultimo barlume dell'Arco sparì, e così anche l'orizzonte. Sul ponte di
comando della \emph{Capetown Maru} la bussola doveva aver iniziato a
ruotare; il capitano suonò la sirena della nave, un rumore brutalmente
assordante, il raglio di uno spazio oltraggiato. Guardai in alto. Le
stelle si mescolavano tra loro provocandoci un capogiro.

«\emph{Ora}» gridò Diane nel frastuono.

Mi sporsi dal corrimano d'acciaio, la sua mano sulla mia, e rovesciammo
la fialetta. Le ceneri fecero una spirale nel vento, colte dalle luci
della nave come fiocchi di neve. Svanirono prima di toccare le
turbolenti acque nere - sparpagliate, volevo credere, nel vuoto che
stavamo invisibilmente attraversando, il punto di sutura, senza oceano,
tra le stelle.

Diane si appoggiò sul mio petto e il latrato della sirena ci risuonò
dentro come una pulsazione finché alla fine si placò.

Poi alzò la testa. «Il cielo» disse.

Le stelle erano nuove e strane.

\separatorefiori

Al mattino salimmo tutti sul ponte: En, i suoi genitori, Ibu Ina, gli
altri passeggeri, persino Jala e un buon numero di membri
dell'equipaggio non più in servizio, per annusare l'aria e sentire il
calore del nuovo mondo.

Poteva essere la Terra, dal colore del cielo e dal calore della luce
solare. Il capo di Port Magellan apparve come una linea irregolare
all'orizzonte, un promontorio roccioso e un paio di spire di fumo
pallido che salivano in verticale e seguivano un vento più alto verso
ovest.

Ibu Ina si unì a noi sul corrimano. En al seguito.

«Sembra così familiare,» disse Ina, «ma dà sensazioni così diverse.»

Ciuffi d'erba attorcigliata seguivano la nostra scia, liberati dalla
terraferma di Equatoria dalle tempeste o dalle maree, enormi foglie a
otto punte si afflosciavano sulla superficie dell'acqua. L'Arco a quel
punto era dietro di noi, non più una porta d'uscita ma una porta di
rientro, un tipo di passaggio del tutto diverso.

Ina esclamò: «È come se una storia fosse finita e un'altra cominciata.»

En non era d'accordo. «No» la contraddisse solennemente, sporgendosi nel
vento come se potesse prevedere in anticipo il futuro. «La storia non
inizia finché non sbarchiamo.»
